\documentclass[conference]{IEEEtran}
\IEEEoverridecommandlockouts
\usepackage{cite}
\usepackage{amsmath,amssymb,amsfonts}
\usepackage{graphicx}
\usepackage{textcomp}
\usepackage{xcolor}
\usepackage{algorithm}
\usepackage{algpseudocode}
\usepackage{booktabs}
\usepackage{tabularx}

\def\BibTeX{{\rm B\kern-.05em{\sc i\kern-.025em b}\kern-.08em
    T\kern-.1667em\lower.7ex\hbox{E}\kern-.125emX}}

\begin{document}

\title{ModelForge AI Studio: A Browser-Native No-Code Platform for Transfer-Learning-Based Edge Image Classification}

\author{
\IEEEauthorblockN{Varsha Dange}
\IEEEauthorblockA{\textit{Department of Artificial Intelligence and Machine Learning} \\
\textit{Vishwakarma Institute of Technology} \\
Pune, India \\
varsha.dange1@vit.edu}
\and
\IEEEauthorblockN{Tanmay Kumbhare}
\IEEEauthorblockA{\textit{Department of Artificial Intelligence and Machine Learning} \\
\textit{Vishwakarma Institute of Technology} \\
Pune, India \\
tanmay.kumbhare24@vit.edu}
\and
\IEEEauthorblockN{Kasturi Rasekar}
\IEEEauthorblockA{\textit{Department of Artificial Intelligence and Machine Learning} \\
\textit{Vishwakarma Institute of Technology} \\
Pune, India \\
rasekar.kasturi24@vit.edu}
\and
\IEEEauthorblockN{Hemantkumar S. Lakhane}
\IEEEauthorblockA{\textit{Department of Artificial Intelligence and Machine Learning} \\
\textit{Vishwakarma Institute of Technology} \\
Pune, India \\
hemantkumar.lakhane24@vit.edu}
\and
\IEEEauthorblockN{Dhanashri Khode}
\IEEEauthorblockA{\textit{Department of Artificial Intelligence and Machine Learning} \\
\textit{Vishwakarma Institute of Technology} \\
Pune, India \\
dhanashri.khode24@vit.edu}
\and
\IEEEauthorblockN{Yash Lakhe}
\IEEEauthorblockA{\textit{Department of Artificial Intelligence and Machine Learning} \\
\textit{Vishwakarma Institute of Technology} \\
Pune, India \\
yash.lakhe24@vit.edu}
}

\maketitle

\begin{abstract}
The democratization of artificial intelligence relies heavily on accessible tooling that bridges the gap between complex neural architectures and non-expert developers. Traditional machine learning workflows demand substantial coding expertise, complex environment configurations, and access to significant cloud computing resources. To address these barriers, we present ModelForge AI Studio, a browser-oriented, no-code machine learning workspace tailored for image classification. Beyond generating classifiers, ModelForge implements an integrated transparency layer---encompassing stage-wise pipeline introspection, embedding-guided dataset assessment, epoch replay, and occlusion-sensitivity heatmaps---to foster pedagogical explainability. Leveraging TensorFlow.js and MobileNet-based transfer learning, the platform facilitates human-verified dataset refinement through dynamic labeling and multi-level data augmentation. By relocating the computational workload to the client edge via WebGL acceleration, the architecture adopts a resource-aware design that supports a privacy-preserving local workflow and reduces cloud dependency for training and inference. This paper details the system architecture, mathematical formulations, and a formal evaluation protocol for ModelForge AI Studio, positioning it as an accessible and interpretable edge AI solution.
\end{abstract}

\begin{IEEEkeywords}
Edge AI, Transfer Learning, No-Code Machine Learning, TensorFlow.js, Explainable AI, MobileNet, Browser-Based Inference.
\end{IEEEkeywords}

\section{Introduction}

The rapid advancement of deep learning has revolutionized computer vision, enabling breakthroughs in object recognition, medical imaging, and autonomous systems. However, a significant barrier to entry persists: training state-of-the-art models traditionally requires advanced programming skills, specialized hardware, and complex software dependencies. This technical overhead limits the participation of domain experts, educators, and hobbyists who lack specialized knowledge in data science but possess valuable domain-specific insights.

Furthermore, conventional machine learning (ML) paradigms are predominantly cloud-centric. Data is transmitted to centralized servers for processing, raising significant privacy concerns and latency bottlenecks. For applications involving sensitive or proprietary image data, routing datasets through external networks is often unacceptable. Edge computing has emerged as a compelling alternative, processing data closer to its source. The modern web browser, equipped with hardware-accelerated graphics processing capabilities (WebGL) and optimized JavaScript engines, represents a ubiquitous and powerful edge execution environment.

We introduce \textit{ModelForge AI Studio}, a decentralized, browser-native no-code platform designed to democratize image classification. Built on TensorFlow.js, the workspace abstracts the complexities of the ML pipeline while maintaining rigorous performance. Users can visually upload datasets, apply dynamic preprocessing and augmentation, and train custom classifiers using transfer learning via a pre-trained MobileNet backbone. Because training images are not uploaded to an application server during the local workflow, the platform supports a privacy-preserving local workflow and reduces cloud dependency for training and inference.

Unlike some existing no-code tools that function as opaque black boxes, ModelForge AI Studio emphasizes an integrated transparency layer. The platform provides occlusion-sensitivity heatmaps, real-time confusion matrices, a comprehensive pipeline inspector, and embedding distance analysis, allowing users to visually interrogate network decisions and understand feature representations. By bridging the gap between no-code accessibility and deep technical insight, this work presents a practical framework for pedagogical and edge ML deployment.

As illustrated in Fig. \ref{fig:architecture}, the ModelForge architecture is fully browser-native and client-side. Data enters through user uploads or a live webcam, and is subsequently processed via a dynamic labeling studio, preprocessing, and multi-level augmentation. Following MobileNet embedding extraction and classifier training, the system integrates evaluation and XAI before culminating in a TFJS model export and customer live-link workflow. This architecture significantly reduces backend dependencies because the primary ML operations run locally in the browser.

\begin{figure}[ht]
\centering
\includegraphics[width=0.95\linewidth]{fig_architecture.png}
\caption{System Architecture of ModelForge AI Studio: Demonstrating the end-to-end browser-based workflow from data collection to model export.}
\label{fig:architecture}
\end{figure}

\section{Literature Review}

The development of \textit{ModelForge AI Studio} is informed by prior research in transfer learning, no-code frameworks, and edge computing.

\subsection{Transfer Learning and Mobile Architectures}
Transfer learning has proven indispensable for domains with constrained data availability \cite{Shao2015}. By utilizing features learned from massive datasets (e.g., ImageNet), networks can rapidly adapt to new tasks. Howard et al. introduced \textit{MobileNets}, which leverage depthwise separable convolutions to drastically reduce model size and computational complexity without severe accuracy degradation \cite{Howard2017}. This efficiency is critical for mobile and edge environments. Dewan et al. further demonstrated that lightweight architectures like MobileNet and DenseNet are highly effective for transfer learning across varied classification tasks \cite{Dewan2023}.

\subsection{Democratization and Interactive Machine Learning}
To make ML accessible, researchers have explored Interactive Machine Learning (IML) and graphical interfaces. Carney et al. presented \textit{Teachable Machine}, a seminal web-based tool that enables rapid model creation via webcam input without coding \cite{Carney2020}. While highly accessible, such tools often obscure the underlying mechanics. More recent efforts, such as the \textit{asanAI} toolkit \cite{Koch2025}, have introduced offline-first capabilities and educational components. ModelForge AI Studio builds upon this foundation by integrating an advanced labeling studio, granular augmentation controls, and deep inspection tools (e.g., embedding analysis) to provide a feature-rich workspace that remains accessible to novices.

\subsection{Edge Inference and Privacy}
The migration of machine learning to the browser was significantly accelerated by the introduction of \textit{TensorFlow.js} \cite{Smilkov2019}. By utilizing WebGL, TensorFlow.js achieves hardware-accelerated execution for inference and training on client hardware. Because training images are not uploaded to an application server during the local workflow, this paradigm supports a privacy-preserving local workflow. Lian et al. explored similar browser-based frameworks for cross-platform federated learning, highlighting the security benefits of local execution \cite{Lian2021}.

\section{System Architecture and Methodology}

ModelForge AI Studio is architected as a modular, client-side application. The workflow is partitioned into sequential stages: data management, feature extraction, classifier training, and explainable evaluation.

The core of the platform is a human-in-the-loop workflow, depicted in Fig. \ref{fig:workflow}. Users can iteratively upload, label, edit, augment, train, evaluate, explain, and export models. This workflow is explicitly designed for non-expert users, while still exposing critical technical feedback mechanisms such as performance metrics and XAI heatmaps.

\begin{figure}[ht]
\centering
\includegraphics[width=0.95\linewidth]{fig_workflow.png}
\caption{Browser-based ML workflow pipeline highlighting human-in-the-loop dataset refinement, augmentation, and inference.}
\label{fig:workflow}
\end{figure}

\subsection{Human-Verified Augmentation Loop}
The platform features a dynamic Labeling Studio where users import images via local folder uploads or live webcam frames. To mitigate the risk of automated dataset corruption, the system implements a human-verified augmentation loop. Image preprocessing and data augmentation can be applied hierarchically: at the global dataset level to address uniform bias, at the class level to correct specific imbalances, or at the individual image level for localized variation. Transformations include resizing, normalization, color jittering, and geometric shifts. Critically, generated variants are presented to the user for explicit review and acceptance, framing augmentation as an active human-in-the-loop refinement process rather than an opaque, automated mechanism.

\subsection{Transfer Learning via MobileNet Embeddings}
Training a deep Convolutional Neural Network (CNN) from scratch in the browser is computationally prohibitive. Therefore, ModelForge AI Studio employs a transfer learning methodology. 

We utilize a frozen, pre-trained MobileNet backbone to act as a universal feature extractor. For a given preprocessed input image tensor $\mathbf{x} \in \mathbb{R}^{224 \times 224 \times 3}$, the MobileNet backbone $\Phi$ extracts a high-dimensional feature embedding $\mathbf{e}$:
\begin{equation}
\mathbf{e} = \Phi(\mathbf{x})
\end{equation}
where $\mathbf{e} \in \mathbb{R}^{1024}$. This dense 1024-dimensional representation captures the hierarchical visual features of the image.

\subsection{Embedding-Guided Dataset Quality Assessment}
By leveraging the extracted embeddings, the platform enables quantitative analysis of dataset diversity and class separability prior to training. For any defined class $C$ containing $|C|$ embeddings, the class centroid $\boldsymbol{\mu}_C$ is calculated as:
\begin{equation}
\boldsymbol{\mu}_C = \frac{1}{|C|} \sum_{\mathbf{e}_i \in C} \mathbf{e}_i
\end{equation}
Intra-class dispersion $D_{intra}(C)$, which indicates sample diversity, is assessed via the mean Euclidean distance to the centroid:
\begin{equation}
D_{intra}(C) = \frac{1}{|C|} \sum_{\mathbf{e}_i \in C} \|\mathbf{e}_i - \boldsymbol{\mu}_C\|_2
\end{equation}
Furthermore, inter-class separability $D_{inter}(C_a, C_b)$ between classes $a$ and $b$ is computed using cosine distance:
\begin{equation}
D_{inter}(C_a, C_b) = 1 - \frac{\boldsymbol{\mu}_{C_a} \cdot \boldsymbol{\mu}_{C_b}}{\|\boldsymbol{\mu}_{C_a}\| \|\boldsymbol{\mu}_{C_b}\|}
\end{equation}
These mathematical indicators guide users in refining their datasets. High intra-class dispersion coupled with low inter-class separability prompts the user to collect more distinct samples or utilize the human-verified augmentation loop to reinforce weak classes.

To maintain computational efficiency, the system utilizes a transfer-learning structure as shown in Fig. \ref{fig:model_architecture}. The base MobileNet feature extractor remains entirely frozen, and only the dense classification head is trained. The exact custom classifier head comprises a Dense layer with 128 units, a Dropout layer with a rate of 0.3, a second Dense layer with 64 units, and a final Softmax output over $K$ classes. This lightweight structure is suitable for rapid browser-side training.

\begin{figure}[ht]
\centering
\includegraphics[width=0.95\linewidth]{fig_model_architecture.png}
\caption{Architecture detailing the frozen MobileNet feature extractor and the trainable dense classification head.}
\label{fig:model_architecture}
\end{figure}

\subsection{Trainable Classification Head}
The 1024-dimensional embeddings are fed into a lightweight, custom classification head instantiated via TensorFlow.js. This trainable sequential network is structured as follows:
\begin{itemize}
    \item Dense Layer (128 units, ReLU activation)
    \item Dropout Layer (rate = 0.3)
    \item Dense Layer (64 units, ReLU activation)
    \item Dense Output Layer ($K$ units, Softmax activation)
\end{itemize}

The final categorical probability distribution $\mathbf{p}$ over $K$ classes is computed using the softmax activation function. For class $i$, given the raw logit $z_i$:
\begin{equation}
P(y=i \mid \mathbf{x}) = \frac{\exp(z_i)}{\sum_{j=1}^{K} \exp(z_j)}
\end{equation}

During the in-browser training phase, the weights of this classification head are optimized using the Adam optimizer with a learning rate of $0.0005$. Training runs for 25 epochs using a dynamic batch size set to $\min(16, \text{number\_of\_samples})$ and a validation split of $0.1$. The objective function minimizes the categorical cross-entropy loss $\mathcal{L}$:
\begin{equation}
\mathcal{L} = - \sum_{i=1}^{K} y_i \log(P(y=i \mid \mathbf{x}))
\end{equation}
where $y_i$ represents the true binary label (one-hot encoded).

\subsection{Stage-Wise Neural Pipeline Introspection}
To foster pedagogical explainability and provide debugging support, the platform exposes an integrated transparency layer:
\begin{itemize}
    \item \textbf{Pipeline Inspector:} The interface allows users to step through the entire inference lifecycle. It visualizes the raw input frame, the 224$\times$224 resized input, the normalized tensor matrix, the intermediate MobileNet 1024D embedding representation, and the final Softmax probability vector.
    \item \textbf{Occlusion Sensitivity Heatmaps:} The system estimates the spatial importance of image regions by systematically masking patches (size 40$\times$40, stride 24) of the input image and measuring the resultant drop in classification confidence. A significant drop indicates high feature importance in that occluded region.
    \item \textbf{Temporal Interpretability through Epoch Replay:} To interpret training dynamics, the system stores and reconstructs model behavior sequentially over epochs. This allows users to inspect how prediction confidence changes dynamically during training, thereby helping to identify unstable convergence or overfitting behavior.
\end{itemize}

\subsection{Evaluation Metrics Formulation}
The platform generates classification reports post-training. The evaluation incorporates precision, recall, and F1-score derived from the elements of the confusion matrix (True Positives $TP$, False Positives $FP$, False Negatives $FN$):
\begin{equation}
\text{Precision} = \frac{TP}{TP + FP}, \quad \text{Recall} = \frac{TP}{TP + FN}
\end{equation}
\begin{equation}
\text{F1-Score} = 2 \times \frac{\text{Precision} \times \text{Recall}}{\text{Precision} + \text{Recall}}
\end{equation}

\section{Implementation and Deployment}

\subsection{Browser Resource-Aware Design}
The application is constructed utilizing web standards (HTML5, CSS3, JavaScript) and TensorFlow.js. Executing machine learning in the browser introduces strict memory and compute constraints. To mitigate this, ModelForge adopts a resource-aware design. The MobileNet feature extractor remains entirely frozen in memory, significantly reducing the computational training cost. Only the lightweight dense classification head is actively optimized. To achieve low latency, all tensor operations are accelerated via the client's WebGL backend. Active tensor disposal and rigorous memory management protocols are enforced to prevent browser crashes during epoch iteration. Performance fundamentally scales with the host device's hardware constraints and WebGL support.

\subsection{Model Export and Integration}
Upon completion of the training cycle, the trained model (comprising the topology JSON and binary weight files) can be exported to the local filesystem. The system further supports a customer live-link workflow, providing standalone integration guidance for deploying the custom classifier in external web environments.

\subsection{Platform Landscape Comparison}
To contextualize ModelForge AI Studio within the broader ecosystem of automated machine learning solutions, Table \ref{tab:platform_scope} and Table \ref{tab:platform_capabilities} present a divided landscape comparison of prominent no-code and AutoML platforms \cite{Roboflow, LandingLens, VertexAI, AzureAutoML, SageMakerCanvas, DataRobot}. Table \ref{tab:platform_scope} specifically contrasts platform scope and deployment characteristics, while Table \ref{tab:platform_capabilities} highlights browser-local and transparency-oriented workflow capabilities.

It is imperative to distinguish between enterprise-grade cloud AutoML systems and edge-oriented pedagogical tools. Cloud platforms such as Google Vertex AI, Azure AutoML, and DataRobot are designed for robust, large-scale MLOps and inherently rely on remote computing infrastructure. ModelForge is not positioned as a replacement for enterprise AutoML/MLOps platforms, but as a lightweight browser-native workflow designed for education, rapid prototyping, and privacy-sensitive image classification executed entirely on the local edge.

\begin{table*}[t]
\centering
\footnotesize
\caption{Platform Scope and Deployment Characteristics}
\label{tab:platform_scope}
\begin{tabularx}{\textwidth}{@{}p{0.20\textwidth}X p{0.16\textwidth} p{0.18\textwidth}@{}}
\toprule
\textbf{Platform} & \textbf{Scope} & \textbf{Mode} & \textbf{Deployment} \\
\midrule
Teachable Machine \cite{Carney2020} & Image/audio/pose teaching & Browser & TFJS export \\
Roboflow \cite{Roboflow} & CV data/training/deploy & Cloud/edge & Cloud/edge \\
LandingLens \cite{LandingLens} & No-code CV & Cloud/edge & Cloud/edge \\
Vertex AI \cite{VertexAI} & Managed AutoML & Cloud & Vertex AI \\
Azure AutoML \cite{AzureAutoML} & Managed AutoML & Cloud & Azure ML \\
SageMaker Canvas \cite{SageMakerCanvas} & No-code ML & Cloud & SageMaker \\
DataRobot \cite{DataRobot} & AutoML/MLOps & Cloud/on-prem & Registry/monitoring \\
\textbf{ModelForge} & Browser image classification & Browser-local & TFJS/live-link \\
\bottomrule
\end{tabularx}
\end{table*}

\begin{table*}[t]
\centering
\footnotesize
\caption{Browser-Local and Transparency-Oriented Capabilities}
\label{tab:platform_capabilities}
\begin{tabularx}{\textwidth}{@{}p{0.20\textwidth}p{0.12\textwidth}p{0.14\textwidth}X X@{}}
\toprule
\textbf{Platform} & \textbf{No-Code CV} & \textbf{Local Training} & \textbf{Dataset Tools} & \textbf{Transparency} \\
\midrule
Teachable Machine \cite{Carney2020} & Yes & Yes & Limited & Limited \\
Roboflow \cite{Roboflow} & Yes & No & Yes & Platform-dependent \\
LandingLens \cite{LandingLens} & Yes & No & Yes & Review tools \\
Vertex AI \cite{VertexAI} & Yes & No & Limited & Eval tools \\
Azure AutoML \cite{AzureAutoML} & Yes & No & Limited & Eval tools \\
SageMaker Canvas \cite{SageMakerCanvas} & Partial & No & Data prep & Eval tools \\
DataRobot \cite{DataRobot} & Partial & No & Data prep & Governance/XAI \\
\textbf{ModelForge} & Yes & Yes & Multi-level & Pipeline/XAI/replay \\
\bottomrule
\end{tabularx}
\end{table*}

\section{Prototype Evaluation Protocol}

To rigorously validate the efficacy of ModelForge AI Studio, we have established a formal evaluation protocol designed to assess computational efficiency, classification robustness, and user interpretability. The final empirical measurements are pending a complete dataset run. The framework dictates that the platform will be evaluated using a structured, user-provided image dataset.

\subsection{Dataset Preparation}
The evaluation data consists of user-provided labeled image classes. During the protocol, images are systematically organized into respective class folders, imported through the browser's native dataset workflow, and visually reviewed within the dynamic labeling studio. To assess the impact of data diversification, the dataset can be optionally augmented at the global dataset level, class level, or individual image level prior to being split into distinct training and validation subsets.

\subsection{Dataset Summary and Class Distribution}
The current execution involves a total of 581 samples distributed across 5 classes. As detailed in Table \ref{tab:class_distribution}, the dataset is notably imbalanced, with the No-damage category forming the majority class. Because of this skew, macro-averaged precision, recall, and F1-score should be strictly reported in the final evaluation to avoid majority-class bias. Furthermore, class-level augmentation is highly relevant for minority classes such as scratch and broken lights to artificially balance the dataset prior to training.

\begin{table}[ht]
\centering
\caption{Dataset Class Distribution}
\label{tab:class_distribution}
\begin{tabular}{lll}
\toprule
\textbf{Class} & \textbf{Samples} & \textbf{Dataset Share} \\
\midrule
scratch & 38 & 6.5\% \\
No-damage & 357 & 61.4\% \\
dent & 103 & 17.7\% \\
broken lights & 36 & 6.2\% \\
broken glass & 47 & 8.1\% \\
\bottomrule
\end{tabular}
\end{table}

\subsection{Embedding Separability Analysis}
Prior to training, the platform-reported pairwise embedding distances were extracted to diagnose dataset quality (Table \ref{tab:embedding_distance}). The No-damage category is comparatively more separable from damage classes, especially broken glass and broken lights. However, some damage categories show low separability, notably dent vs broken lights and scratch vs dent. This indicates potential visual overlap between damage categories and justifies the need for embedding-guided dataset refinement, targeted augmentation, and rigorous confusion-matrix analysis. These represent pre-training dataset diagnostics and do not dictate final model performance.

\begin{table}[ht]
\centering
\caption{Pairwise Embedding Distances}
\label{tab:embedding_distance}
\begin{tabular}{ll}
\toprule
\textbf{Class Pair} & \textbf{Embedding Distance} \\
\midrule
scratch vs No-damage & 18.2\% \\
scratch vs dent & 5.0\% \\
scratch vs broken lights & 8.4\% \\
scratch vs broken glass & 10.4\% \\
No-damage vs dent & 11.5\% \\
No-damage vs broken lights & 19.2\% \\
No-damage vs broken glass & 21.3\% \\
dent vs broken lights & 4.1\% \\
dent vs broken glass & 7.8\% \\
broken lights vs broken glass & 8.4\% \\
\bottomrule
\end{tabular}
\end{table}

\subsection{Experimental Protocol}
The client-side training pipeline executes under a strictly defined experimental setup. The base MobileNet feature extractor remains entirely frozen, extracting 1024-dimensional spatial embeddings for each input sample. Subsequently, the custom dense classification head is trained for 25 epochs. The optimization process utilizes the Adam optimizer with a learning rate configured to $0.0005$. The system dynamically scales the batch size to $\min(16, \text{number\_of\_samples})$ and reserves a strict $0.1$ validation split for unbiased performance estimation. Formal evaluation will be quantified using accuracy, precision, recall, and F1-score, alongside runtime metrics including inference latency and exported model size.

\subsection{Reporting Template}
Table \ref{tab:evaluation_metrics} presents the formalized reporting template that will be populated upon completion of the dataset execution.

\begin{table}[ht]
\centering
\caption{Evaluation Protocol and Reporting Template}
\label{tab:evaluation_metrics}
\begin{tabularx}{\linewidth}{@{}p{0.45\linewidth} p{0.15\linewidth} X@{}}
\toprule
\textbf{Metric} & \textbf{Value to Report} & \textbf{Notes} \\
\midrule
Number of classes & 5 & User-defined \\
Images per class & 38 / 357 / 103 / 36 / 47 & Before augmentation \\
Total images & 581 & Including splits \\
Augmentation setting & \textit{Pending} & Class/image-level available; final setting pending \\
Validation accuracy & \textit{Pending} & Overall percentage \\
Macro precision & \textit{Pending} & Across all classes \\
Macro recall & \textit{Pending} & Across all classes \\
Macro-F1 & \textit{Pending} & Harmonic mean \\
Training time per epoch & \textit{Pending} & WebGL accelerated \\
Inference time per frame & \textit{Pending} & Real-time latency \\
Exported model size & \textit{Pending} & JSON + bin \\
\bottomrule
\end{tabularx}
\end{table}

\subsection{Expected Analysis}
Following the empirical data collection phase, the results will undergo comprehensive analysis. We will systematically evaluate whether multi-level augmentation demonstrably improves model generalization on the validation set. Furthermore, the generated confusion matrix will be analyzed to identify underlying class overlaps and misclassification patterns. From an interpretability standpoint, we will critically examine whether the generated XAI occlusion heatmaps focus on mathematically meaningful and visually salient image regions. Finally, the timing metrics will be reviewed to confirm whether browser-side training and inference remain highly responsive and usable on standard commodity hardware. The confusion matrix and classification metrics will be inserted after the final training and validation run.

\section{Limitations and Threats to Validity}
While ModelForge AI Studio democratizes local machine learning, its efficacy is bound by specific technical limitations. First, classification performance is heavily dependent on user-provided dataset quality and intra-class balance, as the browser lacks resources to train deep architectures from scratch. Second, as a client-side solution, performance varies significantly based on local hardware capabilities and browser-specific WebGL support. Third, the generation of occlusion-sensitivity heatmaps is computationally intensive and incurs higher latency compared to standard forward-pass predictions. Finally, the current architecture strictly supports image classification tasks and cannot facilitate object detection or segmentation workflows. Comprehensive formal user studies are still required to empirically validate pedagogical effectiveness at scale.

\section{Conclusion}

ModelForge AI Studio demonstrates that interpretable machine learning workflows can be effectively decentralized through modern web technologies. By combining the computational efficiency of transfer learning with the client-side execution of TensorFlow.js, the platform delivers a lightweight, no-code environment for edge image classification. The integration of stage-wise introspection, human-verified data curation, and explainable AI mechanisms bridges the gap between novice accessibility and deep technical insight. Crucially, the resource-aware localized execution supports a local-data workflow while reducing cloud dependency for training and inference. Future iterations will focus on gathering empirical metrics to validate the proposed evaluation protocol.

\section*{Acknowledgment}

The authors express gratitude to the Department of Artificial Intelligence and Machine Learning at Vishwakarma Institute of Technology, Pune, for providing the necessary research environment and academic guidance.

\begin{thebibliography}{00}

\bibitem{Howard2017}
A. G. Howard, M. Zhu, B. Chen, D. Kalenichenko, W. Wang, T. Weyand, M. Andreetto, and H. Adam, ``MobileNets: Efficient Convolutional Neural Networks for Mobile Vision Applications,'' \textit{arXiv preprint arXiv:1704.04861}, 2017.

\bibitem{Smilkov2019}
D. Smilkov, N. Thorat, Y. Assogba, C. Yuan, N. Kreeger, P. Yu, K. Zhang, S. Cai, E. Nielson, D. Soergel, S. Bileschi, et al., ``TensorFlow.js: Machine Learning for the Web and Beyond,'' in \textit{Proc. of the 2nd SysML Conference}, Palo Alto, CA, USA, 2019.

\bibitem{Carney2020}
M. Carney, B. Webster, I. Alvarado, K. Phillips, N. Howell, J. Griffith, J. Jongejan, A. Pitaru, and A. Chen, ``Teachable Machine: Approachable Web-Based Tool for Exploring Machine Learning Classification,'' in \textit{Extended Abstracts of the 2020 CHI Conference on Human Factors in Computing Systems (CHI EA '20)}, Honolulu, HI, USA, 2020, pp. 1--8.

\bibitem{Koch2025}
N. Koch and S. Ghiasvand, ``asanAI: In-Browser, No-Code, Offline-First Machine Learning Toolkit,'' \textit{arXiv preprint arXiv:2501.06226}, 2025.

\bibitem{Shao2015}
L. Shao, F. Zhu, and X. Li, ``Transfer Learning for Visual Categorization: A Survey,'' \textit{IEEE Transactions on Neural Networks and Learning Systems}, vol. 26, no. 5, pp. 1019--1034, 2015.

\bibitem{Dewan2023}
J. H. Dewan, H. Jadhav, R. Das, N. Narsale, S. Nambiar, S. D. Thepade, and A. Mhasawade, ``Image Classification by Transfer Learning using Pre-Trained CNN Models,'' in \textit{Proc. Int. Conf. on Recent Adv. in Electr., Electron., Ubiquitous Commun., and Comput. Intell. (RAEEUCCI)}, Dhaka, Bangladesh, 2023, doi: 10.1109/RAEEUCCI57140.2023.10134069.

\bibitem{Lian2021}
Z. Lian, W. Wang, and C. Su, ``WebFed: Cross-platform Federated Learning Framework Based on Web Browser with Local Differential Privacy,'' \textit{arXiv preprint arXiv:2110.11646}, 2021.

\bibitem{Zeiler2014}
M. D. Zeiler and R. Fergus, ``Visualizing and Understanding Convolutional Networks,'' in \textit{European Conference on Computer Vision (ECCV)}, 2014, pp. 818--833.

\bibitem{Roboflow}
Roboflow, ``Roboflow Train and Deploy Documentation,'' [Online]. Available: https://docs.roboflow.com/train and https://docs.roboflow.com/deploy.

\bibitem{LandingLens}
Landing AI, ``LandingLens Platform Documentation,'' [Online]. Available: https://docs.landing.ai/landinglens.

\bibitem{VertexAI}
Google Cloud, ``Vertex AI AutoML Training Overview,'' [Online]. Available: https://docs.cloud.google.com/vertex-ai/docs/training/automl-training-overview.

\bibitem{AzureAutoML}
Microsoft Azure, ``What is automated machine learning (AutoML)?,'' [Online]. Available: https://learn.microsoft.com/en-us/azure/machine-learning/concept-automated-ml.

\bibitem{SageMakerCanvas}
Amazon Web Services, ``Amazon SageMaker Canvas,'' [Online]. Available: https://aws.amazon.com/sagemaker-ai/canvas/.

\bibitem{DataRobot}
DataRobot, ``DataRobot AI Platform Overview,'' [Online]. Available: https://docs.datarobot.com/en/docs/get-started/day0/orientation.html.

\end{thebibliography}

\end{document}
